%----------------------------------------------------------------------------------------
%	PACKAGES AND THEMES
%----------------------------------------------------------------------------------------

\documentclass{beamer}

\mode<presentation> {

% The Beamer class comes with a number of default slide themes
% which change the colors and layouts of slides. Below this is a list
% of all the themes, uncomment each in turn to see what they look like.

%\usetheme{default}
%\usetheme{AnnArbor}
%\usetheme{Antibes}
%\usetheme{Bergen}
%\usetheme{Berkeley}
%\usetheme{Berlin}
%\usetheme{Boadilla}
%\usetheme{CambridgeUS}
%\usetheme{Copenhagen}
%\usetheme{Darmstadt}
%\usetheme{Dresden}
%\usetheme{Frankfurt}
%\usetheme{Goettingen}
%\usetheme{Hannover}
%\usetheme{Ilmenau}
%\usetheme{JuanLesPins}
%\usetheme{Luebeck}
\usetheme{Madrid}
%\usetheme{Malmoe}
%\usetheme{Marburg}
%\usetheme{Montpellier}
%\usetheme{PaloAlto}
%\usetheme{Pittsburgh}
%\usetheme{Rochester}
%\usetheme{Singapore}
%\usetheme{Szeged}
%\usetheme{Warsaw}

% As well as themes, the Beamer class has a number of color themes
% for any slide theme. Uncomment each of these in turn to see how it
% changes the colors of your current slide theme.

%\usecolortheme{albatross}
%\usecolortheme{beaver}
%\usecolortheme{beetle}
%\usecolortheme{crane}
%\usecolortheme{dolphin}
%\usecolortheme{dove}
%\usecolortheme{fly}
%\usecolortheme{lily}
%\usecolortheme{orchid}
%\usecolortheme{rose}
%\usecolortheme{seagull}
%\usecolortheme{seahorse}
%\usecolortheme{whale}
%\usecolortheme{wolverine}

%\setbeamertemplate{footline} % To remove the footer line in all slides uncomment this line
%\setbeamertemplate{footline}[page number] % To replace the footer line in all slides with a simple slide count uncomment this line

%\setbeamertemplate{navigation symbols}{} % To remove the navigation symbols from the bottom of all slides uncomment this line
}

\usepackage{graphicx} % Allows including images
\usepackage{booktabs} % Allows the use of \toprule, \midrule and \bottomrule in tables
\usepackage[ngerman]{babel}
\usepackage{adjustbox}
\usepackage{tikz}
\usetikzlibrary{arrows.meta}
\tikzset{%
  >={Latex[width=2mm,length=2mm]},
  % Specifications for style of nodes:
            base/.style = {rectangle, rounded corners, draw=black,
                           minimum width=4cm, minimum height=1cm,
                           text centered, font=\sffamily},
  activityStarts/.style = {base, fill=blue!30},
       startstop/.style = {base, fill=red!30},
    activityRuns/.style = {base, fill=green!30},
         process/.style = {base, minimum width=2.5cm, fill=orange!15,
                           font=\ttfamily},
}
%----------------------------------------------------------------------------------------
%	TITLE PAGE
%----------------------------------------------------------------------------------------
\title[Viculum]{"VICULUM" \\ Ein virtueller Rundgang durch die Veranstaltungen eines Hochschul-Curriculums } % The short title appears at the bottom of every slide, the full title is only on the title page

\author{Altenberg, Dürr, Milazzo, Röhrle}
\institute[]
{
Hochschule Albstadt-Sigmaringen \\ % Your institution for the title page
\medskip
\textit{Fabian Altenberg (WIN), altenbfa@hs-albsig.de} % Your email address
\newline
\textit{Maik Dürr (TI), duerrmai@hs-albsig.de} % Your email address
\newline
\textit{Domenico Milazzo (TI), milazzdo@hs-albsig.de} % Your email address
\newline
\textit{Prof. Dr. Jörg Röhrle, roehrle@hs-albsig.de} % Your email address
\newline
}
\date{\today} % Date, can be changed to a custom date

\begin{document}

\begin{frame}
\titlepage % Print the title page as the first slide
\end{frame}

\begin{frame}
\frametitle{Überblick}
\begin{enumerate}
\item Die Idee : Konzept und Realisierung eines virtuellen Rundgangs durch die Veranstaltungen eines Modulhandbuchs 
\item Das Konzept: Ein Datenmodell zur vollständigen Erfassung eines Curriculums 
\item Die Realierung - Datenbankzugriff, Virtualisierung, Implementierung der Kommunikationsschicht
\item Zusammenfassung/Kritik/Ausblick
\end{enumerate}
\end{frame}

%----------------------------------------------------------------------------------------
%	PRESENTATION SLIDES
%----------------------------------------------------------------------------------------

%------------------------------------------------

\begin{frame}
\frametitle{Die Idee}
Implementierung eines natürlichen ("spielerischen") Zugangs zu Studieninhalten, quasi gleichsam als Ersatz der zwar korrekten, jedoch "trockenen" Lektüre des Modulhandbuchs ("Curriculums")
Projektion eines Studienverlaufs auf ein mehrstöckiges Gebäude, dessen Etagen die einzelnen Studiensemester und deren Räume wiederum die darin angebotenen Veranstaltungen repräsentieren 
Recherchen ergaben: Bis dato keine vergleichbaren Realisierungen unserer Mitbewerber!\\~\\
\end{frame}


%------------------------------------------------

\begin{frame}
\frametitle{Die Idee}
\begin{itemize}
\item Oracle Datenbank als Speichermedium
\item Unity als Visualisierung Tool der Virtual Reality
\item Dynamischen Aufbau der Szenen um Objekte
\item Bilder \& Video als zusätzliche Medien
\end{itemize}
\end{frame}


%------------------------------------------------

\begin{frame}
\frametitle{Datenbankmodell}
\begin{center}
\includegraphics[width=0.8\textwidth]{pictures/Datenmodell_Viculum.png}
\end{center}
\end{frame}


%------------------------------------------------

\begin{frame}
\frametitle{Datenzugriff}
\begin{adjustbox}{totalheight=\textheight-3\baselineskip}
\begin{columns}[c] % The "c" option specifies centered vertical alignment while the "t" option is used for top vertical alignment

\column{.45\textwidth} % Left column and width

\textbf{Ablauf}
\begin{enumerate}
\item Szenen wird gestartet und holt sich aus der statischen Klasse die Informationen aus der vorherigen Szene.
\item Leeres Objekt wird angestoßen und mit den Szeneninformation gefüllt.
\item Das C\# Script des leeren Objekt wird ausgeführt.
\item Datenbankzugriff auf die Oracle Datenbank durch Select Statement und deren Auswertung.
\item Die Informationen der Auswertung werden im Objekt angezeigt.
\end{enumerate}

\column{.3\textwidth} % Right column and width
\begin{tikzpicture}[node distance=1.5cm,
    every node/.style={fill=white, font=\sffamily}, align=center]
  % Specification of nodes (position, etc.)
  \node (start)             [activityStarts]              	{Szenen start};
  \node (Object)     [process, below of=start]  {Objekt anstoßen};
  \node (Script)      [process, below of=Object]{Script ausführen};
  \node (DB)     [process, below of=Script]   		{Datenbankzugriff};
  \node (end)      [activityRuns, below of=DB]		{Objekt mit Daten befüllt};   
  % Specification of lines between nodes specified above
  % with aditional nodes for description 
  \draw[->]             (start) -- (Object);
  \draw[->]     		(Object) -- (Script);
  \draw[->]      		(Script) -- (DB);
  \draw[->]     		(DB) -- (end);
  \end{tikzpicture}


\end{columns}
\end{adjustbox}
\end{frame}

%------------------------------------------------

\begin{frame}
\Huge{\centerline{Demo}}
\end{frame}
\end{document} 